\chapter{Involución mínima}

\section{Simetría estructural mínima}
\begin{axiom}
Existe una aplicación \(\Iop\) tal que, para todo estado \(x\) dentro del dominio declarado,
\[
\Iop(\Iop(x)) = x.
\]
\end{axiom}

\begin{observation}
La involución mínima fija una regla de retorno por doble aplicación. Su función en esta base no es agotar toda dualidad posible, sino declarar una simetría estructural mínima de inversión.
\end{observation}

\section{Fijación del origen}
\begin{axiom}
El origen canónico \(\subz\) es punto fijo de la involución:
\[
\Iop(\subz)=\subz.
\]
\end{axiom}

\begin{observation}
La fijación de \subz\ permite conservar una referencia estable al atravesar inversión, oposición o retorno. Esta cláusula no extiende todavía su efecto a todos los regímenes posteriores.
\end{observation}

\section{Inversión orientada}
\begin{axiom}
Cuando una magnitud posee orientación declarada, la involución mínima actúa como inversión de orientación.
\end{axiom}

\begin{observation}
Esta cláusula no debe leerse como ontología global. En la Base Axiomática solo fija una forma mínima de oposición legible cuando el sistema ha declarado orientación.
\end{observation}

\section{Cláusula de alcance}
\begin{observation}
La introducción de \Iop\ establece una simetría mínima suficiente para distinguir inversión, retorno y oposición sin ambigüedad inmediata. No valida por sí sola toda arquitectura ulterior de dualidad, cambio de régimen o colapso estructural.
\end{observation}

